% BHU PYQ -- Manually transcribed & error-corrected
% Paper: BCH-111 : Financial Accounting-I
% Exam: B.Com. (Hons.) Semester I Examination, 2022-23

\documentclass[11pt,a4paper]{article}
\usepackage[a4paper,top=2.5cm,bottom=2.5cm,left=2.8cm,right=2.8cm,headheight=14pt]{geometry}
\usepackage{amsmath,amssymb}
\usepackage[T1]{fontenc}
\usepackage[utf8]{inputenc}
\usepackage{lmodern,microtype,fancyhdr,enumitem,hyperref,lastpage,parskip}
\usepackage{booktabs}

\hypersetup{
    colorlinks=true,
    urlcolor=black,
    linkcolor=black,
    pdftitle={BCH-111: Financial Accounting-I -- BHU B.Com. (Hons.) Sem I 2022-23}
}

\pagestyle{fancy}
\fancyhf{}
\renewcommand{\headrulewidth}{0.4pt}
\renewcommand{\footrulewidth}{0.4pt}
\fancyhead[L]{\small   \textit{Banaras Hindu University}}
\fancyhead[R]{\small   \textit{BCH-111: Financial Accounting-I}}
\fancyfoot[C]{\small Page  \thepage\ of \pageref{LastPage}}


\newlist{parts}{enumerate}{2}
\setlist[parts,1]{label=(\alph*),leftmargin=*,itemsep=5pt,topsep=3pt}
\setlist[parts,2]{label=(\roman*),leftmargin=*,itemsep=3pt,topsep=2pt}

\newcommand{\pts}[1]{\hfill   \textit{[#1]}}


\begin{document}


\begin{center}
  {\large\bfseries BANARAS HINDU UNIVERSITY}\\[2pt]
  {
\normalsize B.Com. (Hons.) --- Semester I Examination, 2022-23}\\[6pt]
  \rule{0.8\linewidth}{0.5pt}\\[6pt]
  {\large\bfseries Paper No. BCH-111: Financial Accounting-I}\\[4pt]
  {
\normalsize Time: 3 Hours\hfill Maximum Marks: 70}
\end{center}

\rule{\linewidth}{0.4pt}

\smallskip



\noindent   \textit{Instructions: Answer all questions. All questions carry equal marks (14 marks each).}

\bigskip




\noindent   \textbf{Question 1.}
Journalise the following transactions in the Journal of S Sons:\pts{14}

\begin{itemize}[itemsep=2pt,topsep=2pt]
  \item[2022, April 1:] S Sons started his business with Capital of Rs. 5,00,000 and Furniture worth Rs. 25,000.
  \item[April 2:] Bought goods of list price of Rs. 80,000 at 10\% trade discount from B Sons.
  \item[April 3:] Opened a current account in Bank with Rs. 1,00,000.
  \item[April 5:] Purchased goods worth Rs. 1,00,000 from Khanna Brothers at a trade discount of 20\% and cash discount of 5\%.
  \item[April 10:] Purchased machinery for Rs. 20,000 and spent Rs. 500 on its carriage and Rs. 300 on its installation.
  \item[April 12:] Sold goods worth Rs. 40,000 to A Sons.
  \item[April 14:] Paid Rs. 3,500 for Life Insurance Premium of the owner.
  \item[April 15:] Returned goods to Khanna Brothers Rs. 10,000.
  \item[April 18:] A Sons became insolvent and only 50 paise in a Rupee were recovered from him.
  \item[April 22:] Withdrew from bank Rs. 4,000 for office use and Rs. 5,000 for personal use.
  \item[April 30:] Salary of Rs. 10,000 is unpaid.
\end{itemize}

\centerline{   \textbf{OR}}

Explain any two of the following:\pts{14}

\begin{parts}
  \item Capital Expenditure
  \item Principle of Conservatism
  \item Principle of Full Disclosure
  \item Separate Entity Concept
\end{parts}

\medskip\rule{\linewidth}{0.2pt}


\newpage




\noindent   \textbf{Question 2.}
From the following Trial Balance and additional information, prepare the Trading and Profit and Loss Account and Balance Sheet for the year ended 31st March, 2022 of S Sons:

\begin{center}
  \small
  
\begin{tabular}{lrr}
     oprule
   \textbf{Particulars} &   \textbf{Debit Balance (Rs.)} &   \textbf{Credit Balance (Rs.)} \\
    \midrule
    Capital Account & & 16,00,000 \\
    Sundry Creditors & & 7,20,000 \\
    Sundry Debtors & 14,50,000 & \\
    Sales & & 28,85,500 \\
    Purchases & 12,69,000 & \\
    General Expenses & 1,50,000 & \\
    Carriage Inward & 60,000 & \\
    Carriage Outward & 90,000 & \\
    Wages & 1,00,000 & \\
    Discount & 30,000 & 49,500 \\
    Bad Debts & 20,000 & \\
    Land and Building & 10,40,000 & \\
    Printing and Stationery & 25,000 & \\
    Furniture & 1,90,000 & \\
    Advertisement & 50,000 & \\
    Insurance & 40,000 & \\
    Postage and Telephone & 40,500 & \\
    Salaries & 2,60,000 & \\
    Rates and Taxes & 20,000 & \\
    Drawings & 20,000 & \\
    Stock on 1-4-2021 & 2,50,000 & \\
    Cash at Bank & 1,30,000 & \\
    Cash in Hand & 20,500 & \\
    \midrule
   \textbf{Total} &   \textbf{52,55,000} &   \textbf{52,55,000} \\
    ottomrule
  \end{tabular}
\end{center}




\noindent   \textbf{Adjustments:}

\begin{parts}
  \item Stock on 31st March, 2022 was valued at Rs. 8,50,000.
  \item A provision of 8\% for Bad and Doubtful Debts is to be created.
  \item Depreciate Furniture by 15\% and Building by 12\%.
  \item The owner of the firm had withdrawn goods worth Rs. 15,000 during the year for personal use.
  \item Purchases include purchase of Furniture worth Rs. 45,000.
\end{parts}
\pts{14}

\centerline{   \textbf{OR}}

Explain the following:\pts{14}

\begin{parts}
  \item Depreciation
  \item Rectification of Errors
\end{parts}

\medskip\rule{\linewidth}{0.2pt}


\newpage




\noindent   \textbf{Question 3.}
A and B are partners in a firm sharing profit and loss in the ratio of 60\% and 40\%. On 1st April, 2022, the position of the business was as follows:

\begin{center}
  \small
  
\begin{tabular}{lr|lr}
     oprule
   \textbf{Liabilities} &   \textbf{Amount (Rs.)} &   \textbf{Assets} &   \textbf{Amount (Rs.)} \\
    \midrule
    Creditors & 15,000 & Cash at Bank & 7,000 \\
    Loan & 20,000 & Debtors & 18,000 \\
    A's Capital & 30,000 & Stock & 20,000 \\
    B's Capital & 25,000 & Land \& Building & 30,000 \\
              &        & Plant & 15,000 \\
    \midrule
   \textbf{Total} &   \textbf{90,000} &   \textbf{Total} &   \textbf{90,000} \\
    ottomrule
  \end{tabular}
\end{center}

On 1st April, 2022, C is admitted as a partner for one-fourth share under the following terms:

\begin{parts}
  \item C is to pay Rs. 20,000 as capital and Rs. 10,000 as goodwill for his share in the firm.
  \item A revaluation of the assets of the firm will be made by reducing Plant by 20\% and Stock by 10\%.
  \item Provision for bad and doubtful debts is to be raised to 10\% on debtors.
  \item The profit sharing ratio of A and B is to remain unaltered.
\end{parts}
You are required to pass necessary journal entries and prepare the Balance Sheet of the new firm after the admission of C.\pts{14}

\centerline{   \textbf{OR}}

Write notes on any two of the following:\pts{14}

\begin{parts}
  \item Accounting for Goodwill when a new partner is admitted
  \item Fixed and Fluctuating Capital Accounts
  \item Admission of a new partner and profit sharing ratio
\end{parts}

\medskip\rule{\linewidth}{0.2pt}


\newpage




\noindent   \textbf{Question 4.}
P, Q and R are partners sharing profit in the ratio of 4:4:2 respectively. On 31st March, 2022 their Balance Sheet stood as follows:

\begin{center}
  \small
  
\begin{tabular}{lr|lr}
     oprule
   \textbf{Liabilities} &   \textbf{Amount (Rs.)} &   \textbf{Assets} &   \textbf{Amount (Rs.)} \\
    \midrule
    Creditors & 3,12,000 & Plant \& Machinery & 8,60,000 \\
    Bank Overdraft & 48,000 & Land \& Building & 3,28,000 \\
    P's Capital & 8,00,000 & Debtors & 4,70,000 \\
    Q's Capital & 8,00,000 & Stock & 7,02,000 \\
    R's Capital & 4,00,000 & & \\
    \midrule
   \textbf{Total} &   \textbf{23,60,000} &   \textbf{Total} &   \textbf{23,60,000} \\
    ottomrule
  \end{tabular}
\end{center}

On the date of the Balance Sheet, R retired and it was agreed that:

\begin{parts}
  \item The value of the goodwill of the firm is to be fixed at Rs. 5,00,000.
  \item P and Q will share profit in future in the ratio of 3:2.
  \item The amount due to R will be paid immediately and for this purpose P and Q will bring in sufficient cash.
\end{parts}
You are required to pass necessary journal entries and prepare the Balance Sheet of the new firm.\pts{14}

\centerline{   \textbf{OR}}

Explain as to how amount payable to a deceased partner is determined.\pts{14}

\medskip\rule{\linewidth}{0.2pt}


\newpage




\noindent   \textbf{Question 5.}
Explain the Garner vs. Murray rule assuming that A, B and C are equal partners, C has become insolvent, and their Balance Sheet is as follows:

\begin{center}
  \small
  
\begin{tabular}{lr|lr}
     oprule
   \textbf{Liabilities} &   \textbf{Amount (Rs.)} &   \textbf{Assets} &   \textbf{Amount (Rs.)} \\
    \midrule
    A's Capital & 40,00,000 & Cash & 30,00,000 \\
    B's Capital & 12,00,000 & C's Capital & 4,00,000 \\
                &           & Loss on Realisation & 18,00,000 \\
    \midrule
   \textbf{Total} &   \textbf{52,00,000} &   \textbf{Total} &   \textbf{52,00,000} \\
    ottomrule
  \end{tabular}
\end{center}
Close the books of the firm applying the Garner vs. Murray rule.\pts{14}

\centerline{   \textbf{OR}}

Write notes on any one of the following:\pts{14}

\begin{parts}
  \item Dissolution of firm
  \item Piecemeal Distribution on Dissolution
\end{parts}

\vfill
\rule{\linewidth}{0.4pt}

\begin{center}\small
\href{https://bhu-exam.vercel.app/}{Click here to visit our website}\\[4pt]
\href{https://me-aryan.vercel.app/}{Click here to visit our contributor}\\[4pt]
\href{https://quashan-me.vercel.app/}{Click here to visit our contributor}
\end{center}

\end{document}
